\section{Finite computations and split expressions}\label{sec:algebra}

Our first task is to express a finite computation on a sequence of words
using tests on the two sides of selected cuts. The essential requirement is
universal soundness: any independently accepted prefix and suffix must describe
the same deterministic update. We first prove the expression identity that
uses this requirement, and then give two ways to arrange the updates, by
affine correction and by predicted products.

Every test reads the original alphabet. Auxiliary finite sets index unions,
intersections, and state queries; they do not introduce input letters.

\subsection{Words, monoids, and graph tests}

If \(\Sigma=\varnothing\), the only languages are \(0\) and \(1\), and
Theorem~\ref{thm:main} is immediate. Assume henceforth that \(\Sigma\ne\varnothing\).
Let \(M\) be a finite monoid and let \(T:\Sigma^*\to M\) be a morphism.
Positions are integer gaps; the factor on \([a,b)\) has length \(b-a\) and
monoid value \(T_{a,b}\). Composition is in reading order:
\[
                  T_{a,b}T_{b,c}=T_{a,c}.
\]
A right action on a vector space is denoted \(v\mapsto vT_{a,b}\).
The vector space with basis indexed by \(M\) has the faithful action
\(e_m d=e_{md}\); the image of \(e_{1_M}\) identifies \(d\).
All vector spaces used below are finite and are over \(\mathbf F_2\).

For a DFA, take \(M\) to be its monoid of state transformations induced by words,
including the identity. The product of a word determines acceptance. It is
therefore enough to construct uniformly bounded-height tests for each value of
\(T\). The theorem will follow by a finite accepting union.

\begin{remark}[A DFA-state representation]\label{rem:dfa-basis}
One may instead begin with the basis indexed by DFA states and query only
the image of the initial-state vector to test acceptance. That image need not
identify the monoid product. The main proof uses the basis indexed by $M$;
the alternative and its final acceptance query are described at the end
of Section~\ref{sec:compilation}. In either case interval predictions and
the later work registers still take values in $M$.
\end{remark}

A \emph{prefix code} \(E\) is a set of nonempty words, no one a proper prefix of
another. Every word in \(E^*\) has a unique factorization into members of \(E\):
compare first factors, one of which would otherwise properly prefix the other,
and continue inductively. Members of \(E\) will be called \emph{episodes}.
For \(D\subseteq E\), let each \(e\in D\) have a total update \(G_e:S\to S\) on a
finite state set. For a sequence \(w=e_1\cdots e_r\), let \(G_w\) be their
composition in reading order. The empty sequence has the identity update.
The \emph{graph test} from \(x\) to \(y\) on \(D^*\) is
\[
                     \{w\in D^*:G_w(x)=y\}.
\]
Totality and determinism refer to the word itself, not to a chosen factorization
of a test expression or a successful metadata guess.

Intersections stand for
\(\neg(\neg P+\neg Q)\), and other Boolean operations have analogous expansions.
They do not increase height. A fixed finite language is a finite union of
concatenations of letters, with \(1\) for the empty word, so has height zero.
We use this observation only when the \emph{whole argument} of the test has
bounded length.

\subsection{The split identity}

\begin{lemma}[Split identity]\label{lem:split}
Let \(D'\subseteq D\subseteq E\), with the same deterministic updates on the
common episodes. Suppose languages \(P_x,Q_y\), indexed by \(x,y\in S\), satisfy:
\begin{enumerate}[label=\textup{(\roman*)}]
\item If \(w\in D^*\) and \(uv\) is a prefix of \(w\), with
\(u\in P_x\) and \(v\in Q_y\), then \(w\) has a first episode \(e\),
\(uv=e\), and \(G_e(x)=y\).
\item For every \(e\in D\setminus D'\) and \(x\in S\), there is a split
\(e=uv\) with \(u\in P_x\) and \(v\in Q_{G_e(x)}\).
\end{enumerate}
Write \(W_{x,y}\) for the exact graph test on \((D')^*\). Then the exact graph
test on \(D^*\) is
\begin{equation}\label{eq:split}
\begin{split}
D^*\cap\Bigl(
 W_{x,y}\ \cup\
 &\bigl(\,\bigcup_a W_{x,a}P_a\,\bigr)
 \bigl(\,\bigcup_{b,c}Q_bW_{b,c}P_c\,\bigr)^*\\[-2pt]
 &\hspace{35mm}\cdot
 \bigl(\,\bigcup_r Q_rW_{r,y}\,\bigr)\Bigr).
\end{split}
\end{equation}
In particular, if \(D,P_x,Q_y\) have heights at most \(1,0,1\), respectively,
and all \(W_{x,y}\) have height at most \(H\), the new graph tests have height
at most \(\max\{2,H+1\}\).
\end{lemma}

\begin{proof}
Consider a word accepted by the right side of \eqref{eq:split}. Its membership
in \(D^*\) fixes its actual episode parsing. In a factorization through the
nontrivial branch, the initial \(W\)-factor is a sequence of \(D'\)-episodes.
Prefix-code uniqueness makes these actual initial episodes. The adjacent
\(P,Q\) factors now begin at a true boundary; condition~(i), applied to the
remaining word in \(D^*\), makes them exactly the next episode and gives its
indexed update. The next \(W\)-factor consumes actual skipped episodes.
Inducting through the finite number of starred factors proves the claimed
whole update. The \(W_{x,y}\) branch handles words without an explicit split.

Conversely, for a genuine episode sequence, split every episode outside \(D'\)
using~(ii) and collect consecutive skipped episodes into \(W\)-factors.
The deterministic computation supplies the intervening state indices.
If every episode is skipped, use \(W_{x,y}\). This also handles \(\epsilon\).

For the height statement, \(D^*\) has height at most two. The starred middle
factor has height at most \(1+\max\{1,H\}\). All remaining operations and
factors fit the asserted maximum.
\end{proof}

Condition~(i) quantifies over \emph{every} pair of independent successful
witnesses. Equality of guessed parameters is never part of the ambient
meaning of a split. It must follow from the actual tests, or its failure must
be harmless for the declared update.

When \(D'=\varnothing\), the bottom graph tests are \(1\) for equal states and
\(0\) for unequal states. Thus the recursion has height-zero initial tests.

\subsection{Affine corrections and linear recovery}

\begin{lemma}[Affine recovery]\label{lem:affine-recovery}
Suppose the deterministic update on each episode is
\(v\mapsto vA_e+b_e\) on a fixed finite vector space. Exact graph tests for
the composed affine update on all vectors give, by finite Boolean operations,
the graph tests for the composition of the \(A_e\).
\end{lemma}

\begin{proof}
The composite of \(v\mapsto vA+b\) followed by \(v\mapsto vC+d\) is
\(v\mapsto vAC+bC+d\). Therefore the linear part of the composed update
\(g_w\) is \(g_w(v)-g_w(0)\), with the desired product of linear parts.
If \(R_{x,z}\) are its graph tests, then
\[
          \bigcup_z\bigl(R_{x,y+z}\cap R_{0,z}\bigr)
\]
is exactly the test for linear output \(y\) on input \(x\).
Both tests concern the same word and hence the same map.
\end{proof}

An affine map \(f(v)=vA+b\) has the homogeneous linear lift
\begin{equation}\label{eq:homogeneous}
            (v,c)\longmapsto(vA+cb,c)
\end{equation}
on \(V\oplus\mathbf F_2\). These lifts compose and recover \(f\) on inputs
\((v,1)\). More generally, any map of a finite set has a unique linear
extension sending each basis vector to the basis vector of its image.
Neither device requires invertibility.

\subsection{Prediction trials}

Fix a representation of \(M\) on a finite vector space \(U\), and put
\(J=|U|^2\). A trial has endpoints and checkpoints
\[
                 B=z_0<z_1<\cdots<z_J<z_{J+1}=C.
\]
Its block \(j\) is an ordered set of potential cuts strictly between
\(z_{j-1}\) and \(z_j\). Enumerate all pairs of vectors as
\((x_j,y_j)\), \(1\le j\le J\). A prediction is a tuple
\((p_1,\ldots,p_{J+1})\in M^{J+1}\), and \(T_*=p_1\cdots p_{J+1}\).
Write \(T=T_{B,C}\).

\begin{definition}\label{def:eligibility}
Block \(j\) is \emph{ideally eligible} if
\[
 x_jT_*=y_j,\qquad
 T_{z_{i-1},z_i}=p_i\ (1\le i<j),\qquad
 T_{z_j,C}=p_{j+1}\cdots p_{J+1}.
\]
A trial is \emph{perfectly predicted} if every interval product is its
corresponding \(p_i\).
\end{definition}

\begin{lemma}[Prediction trial]\label{lem:prediction}
There is a deterministic shift \(\beta\in U\), defined from the actual trial
word and its prediction, such that every ideally eligible pair satisfies
\(y_j=x_jT+\beta\). A perfectly predicted trial has an eligible block for
every possible input vector.
\end{lemma}

\begin{proof}
Suppose \(h<j\) are eligible. The prefix conditions for \(j\) determine every
actual interval through \(z_h\), and the suffix condition for \(h\) determines
the product from \(z_h\) to \(C\). Hence
\[
 T=T_{B,z_h}T_{z_h,C}
  =(p_1\cdots p_h)(p_{h+1}\cdots p_{J+1})=T_*.
\]
This uses multiplication, without cancellation.

If \(T=T_*\), every eligible pair fits \(\beta=0\). If \(T\ne T_*\), at most
one block is eligible; choose its forced shift \(y_j-x_jT\), or zero if none
is eligible. These rules define \(\beta\) on every trial.
For a perfect prediction, the block for \((x,xT)\) is eligible for each \(x\).
\end{proof}

The conclusion for two eligible blocks is equality of the \emph{whole}
products \(T=T_*\). It does not assert that the individual factors in
the prediction tuple are correct. Perfect prediction is the separate,
stronger condition used to obtain an eligible block for every input.

We replace the action on the trial by \(v\mapsto vT+\beta\), implemented by
adding \(\beta\) at its end. At a cut in block \(j\), the prefix can verify
arrival at \(B\) with vector \(x_j\), the predicted transfer, and every preceding
interval condition. The suffix can verify the product from \(z_j\) to \(C\)
and continue from vector \(y_j\) at \(C\). If these checks agree on the trial
data, Lemma~\ref{lem:prediction} proves their transfer without either side
knowing the entire current trial shift.

Whenever a computation passes through an entire other trial, its shift is
computed using \emph{ideal} eligibility from Definition~\ref{def:eligibility}.
Tag failures, decoding failures, and the choice of a particular split never
affect that shift. This convention will keep all episode updates deterministic.

